IEEE Biocoating Internship With Dr. August 

Omar Ahmadzada
Science & Engineering Manalapan High School Englishtown, NJ 


Shreena Desai
Science & Engineering
Manalapan High School
Englishtown, NJ

Abstract
Implant biofilm infections can change the local environment around an implant, including the pH and amount of available oxygen. Our project looks at whether these changes can affect the behavior of a hydrogel coating and its release of a model drug. We are interested in a coating that could be useful during the period a patient is in the hospital, since infection is highest right after surgery.
We plan to make a PVA-borate hydrogel coating on a stainless-steel 316 washer surface and use FD&C Blue No. 1 dye as a model payload. Before testing environmental responses, we will make several PVA and borate formulations and compare their basic properties. These include how well the coating stays on the surface, how much it swells in solution, how stable it remains, and how consistently it can be made. We will then measure how much dye is released over time at approximately 37 °C in a solution.
After identifying a desired formulation, we will investigate the effect of pH and dissolved oxygen separately. The pH experiments will compare an infection-like acidic condition. We will also test an oxygen-sensitive metalloporphyrin under controlled oxygen conditions before considering its incorporation into the coating. Release measurements will be collected using spectrophotometry or image analysis.
Our hypothesis is that increasing the stability of the PVA-borate network will reduce the initial rate of dye diffusion, while an acidic and oxygen-reduced environment will increase release compared with the control condition. We also expect the response to depend on the PVA and borate concentrations and the thickness of the coating.
This project will not be used to diagnose infection or establish that an infection is present.

Index Terms 
PVA-borate hydrogel, responsive hydrogel, controlled drug release, implant coating,stainless steel 316, model drug delivery, pH-responsive release, oxygen-responsive materials, diffusion, swelling, biomaterials, infection-like conditions.
